\subsection{随机变量的数字特征}

\subsubsection{数学期望}

\begin{enumerate}
  \item 设随机变量 $X$ 的分布函数为
  \begin{equation*}
    F(x)=
    \begin{cases}
      0,   & x<-2\\
      0.3, & -2\le x<1\\
      0.8, & 1\le x<2\\
      1,   & x\ge 2
    \end{cases},
  \end{equation*}
  且 $Y=X^{2}-1$, 则 $E(XY)=(\quad)$
  \begin{choices}
    \item 1.5
    \item 0.3
    \item $-0.6$
    \item 0.6
  \end{choices}

  \item 设随机变量 $X$ 和 $Y$ 相互独立, 二者的概率密度分别为
  \begin{equation*}
    f_X(x)=
    \begin{cases}
      \dfrac{8}{x^{3}}, & x>2\\
      0,                & x\le 2
    \end{cases},
    \qquad
    f_Y(y)=
    \begin{cases}
      2y, & 0<y<1\\
      0,  & \text{其他}
    \end{cases}
  \end{equation*}
  则 $E(XY)=(\quad)$
  \begin{choices}
    \item $\dfrac{2}{3}$
    \item $\dfrac{4}{3}$
    \item $\dfrac{8}{3}$
    \item $\dfrac{10}{3}$
  \end{choices}

  \item 对随机变量 $X$ 和 $Y$, 下面说法错误的是 $(\quad)$
  \begin{choices}
    \item $E(C)=C$
    \item $E(X+Y)=E(X)+E(Y)$
    \item $E(aX+bY)=aE(X)+bE(Y)$
    \item $E(XY)=E(X)E(Y)$
  \end{choices}

  \item 设随机变量 $X$ 的概率密度为
  \begin{equation*}
    f(x)=
    \begin{cases}
      e^{-x}, & x>0\\
      0,      & x\le 0
    \end{cases},
  \end{equation*}
  则 $Y=-2X+1$ 的数学期望为 $(\quad)$
  \begin{choices}
    \item 0
    \item $-1$
    \item 1
    \item $-2$
  \end{choices}

  \item 设随机变量 $(X,Y)$ 的概率密度为
  \begin{equation*}
    f(x,y)=
    \begin{cases}
      12y^{2}, & 0\le y\le x\le 1\\
      0,       & \text{其他}
    \end{cases}
  \end{equation*}
  则 $E(Y)=(\quad)$.
  \begin{choices}
    \item $\dfrac{1}{2}$
    \item $\dfrac{2}{3}$
    \item $\dfrac{3}{5}$
    \item $\dfrac{4}{5}$
  \end{choices}

  \item 设随机变量 $X$ 服从参数为 $1$ 的泊松分布, 则 $P\{X=E(X^{2})\}=(\quad)$
  \begin{choices}
    \item $e^{-1}$
    \item $e^{-1}/2$
    \item $e$
    \item $e/2$
  \end{choices}

  \item 设随机变量 $X$ 和 $Y$ 的关系为 $Y=2X+2$, 若 $X\sim U(0,4)$, 则 $E(Y)=(\quad)$
  \begin{choices}
    \item 4
    \item 6
    \item 8
    \item 10
  \end{choices}

  \item 设随机变量 $X\sim U(0,1)$, 则 $E[X-E(X)]^{2}=(\quad)$
  \begin{choices}
    \item 0
    \item 1
    \item $\dfrac{1}{2}$
    \item $\dfrac{1}{12}$
  \end{choices}

  \item 某射手射击所得环数 $X$ 的分布律如下:
  \begin{center}
    \begin{tabular}{|c|c|c|c|c|}
      \hline
      $X$ & 7   & 8   & 9   & 10  \\
      \hline
      $P$ & $x$ & 0.1 & 0.3 & $y$ \\
      \hline
    \end{tabular}
  \end{center}
  已知 $X$ 的期望 $E(X)=8.9$, 则 $y$ 的值为 $(\quad)$
  \begin{choices}
    \item 0.4
    \item 0.6
    \item 0.7
    \item 0.9
  \end{choices}

  \item 设 $X$ 与 $Y$ 相互独立, 且 $X\sim N(\mu_1,\sigma_1^{2})$, $Y\sim N(\mu_2,\sigma_2^{2})$, 令 $Z=aX+bY+c$, 则 $E(Z)=(\quad)$
  \begin{choices}
    \item $\mu_1+\mu_2$
    \item $\mu_1+\mu_2+c$
    \item $a\mu_1+b\mu_2$
    \item $a\mu_1+b\mu_2+c$
  \end{choices}
\end{enumerate}


\subsubsection{方差}



\subsubsection{随机变量的其他数字特征}